\section*{Parte 4}

En esta sección se analizó una señal FM. Una señal FM viene dada por la siguiente ecuación:

\begin{equation*}
S_{FM}(t)= A cos(2\pi f t + \betha cos(2\pi f t))
\end{equation*}

Luego la misma, al escribirse como forma exponencial, esta compuesta por exponenciales de varias frecuencias a medida que el coseno interno suma o resta la en la fase del coseno de la señal. Aquí se trabajó con una señal con $f_m=1,25MHz$ y $f_c=1,25MHz$. Por tal motivo, al observar la figura \ref{fig:FM}, se puede ver un barrido de frecuencias que van desde $f=0Hz$ hasta $f=f_m+f_c = 2,5MHz$.\par


\begin{figure}[H]
    \centering
    \includegraphics[width=0.8\textwidth]{FOTOS_TP6/PARTE 4/ESTASII.jpeg}
    \caption{Espectro de una señal de fourier}
    \label{fig:FM}
\end{figure}

\par

Finalmente, al comparar con lo visto en el caso de FM, se puede observar lo siguiente. AM, es modulación en amplitud, por lo que la señal se transmite en la amplitud de la señal y se observan solo tres frecuencias. La frecuencias resultante de la multiplicación de la portadora y la modulante más la suma de la frecuencia de la portadora. Por otra parte, en el caso de la FM, como es modulación en frecuencia, se observa al analizar el espectro, que la frecuencia de la misma varía demasiado en comparación a la AM. 